% Capítulo 4 (Proposta) — Slide — Formulação matemática do AIPE
\begin{frame}{Formulação matemática do AIPE}
  \footnotesize
  \textbf{\textcolor{darkblue}{O AIPE formula a reposição como um problema de seleção contextual de política por série loja-produto.}}
  \vspace{0.15em}

  \begin{itemize}\setlength\itemsep{0pt}\setlength\parskip{0pt}
    \item $i \in \mathcal{I}$: série loja-produto;\quad $\pi \in \mathcal{P}$: política candidata
    \item $\mathbf{x}_i$: características operacionais;\quad $g_i$: perfil (POD) de $i$
    \item $\overline{\mathrm{CTI}}_{i,\pi}$, $\overline{\mathrm{NS}}_{i,\pi}$: médias simuladas;\quad $\alpha_{\min}=0{,}70$: nível mínimo de serviço
    \item $z_{i,\pi}\in\{0,1\}$: variável de decisão -- $1$ se $\pi$ é escolhida para $i$; $0$, caso contrário
  \end{itemize}
  \vspace{0.15em}

  \textbf{Formulação por série $i$:}
  \[
    \min_{z_{i,\pi}} \sum_{\pi \in \mathcal{P}} z_{i,\pi}\,\overline{\mathrm{CTI}}_{i,\pi}
  \]
  \vspace{-0.7em}
  sujeito a\;\;
  $\sum_{\pi} z_{i,\pi}=1$,\;\;
  $\sum_{\pi} z_{i,\pi}\,\overline{\mathrm{NS}}_{i,\pi}\geq\alpha_{\min}$,\;\;
  $z_{i,\pi}\in\{0,1\}$
  \vspace{0.18em}

  \textbf{Forma equivalente (gera o rótulo):}\\[1pt]
  $y_i = \arg\min_{\pi \in \mathcal{P}:\,\overline{\mathrm{NS}}_{i,\pi} \geq \alpha_{\min}} \overline{\mathrm{CTI}}_{i,\pi}$
  \qquad $f_{\theta}(\mathbf{x}_i) \rightarrow \hat{\pi}_i \approx y_i$
  \vspace{0.18em}

  \textcolor{warnred}{\textbf{Cuidado conceitual:}} O AIPE não substitui a
  política de reposição: ele escolhe qual política aplicar em cada série
  loja-produto.
  \note{
    Esta é a formulação formal do problema que o AIPE resolve. Para cada
    série loja-produto $i$, queremos escolher uma política $\pi$ -- a
    variável binária $z_{i,\pi}$ -- que minimize o CTI médio, respeitando
    duas restrições: exatamente uma política deve ser escolhida, e essa
    política deve atingir o nível de serviço mínimo de 0,70. Como o
    problema é separável por série e o portfólio é pequeno (12 políticas),
    essa formulação tem solução direta por enumeração: a forma equivalente
    com $\arg\min$ é exatamente isso, e é o que usamos para gerar o rótulo
    $y_i$ no benchmark simulado. O PSE, $f_\theta$, aprende uma função das
    características $\mathbf{x}_i$ que aproxima esse rótulo para séries
    novas, produzindo a recomendação $\hat{\pi}_i$. É importante reforçar:
    o AIPE escolhe qual política usar -- as decisões de quando e quanto
    pedir são geradas pela política escolhida, não pelo AIPE.
  }
\end{frame}
